\chapter{总结与展望}

\section{研究总结}

HBV感染性病毒颗粒的形成依赖成熟核衣壳与包膜蛋白之间的特异性识别，其中L-HBsAg preS区域与HBc核衣壳的相互作用是病毒包膜化过程中的关键环节。近期结构研究提出preS 87--114片段可作为基质样结构域识别HBc表面沟槽，并指出R102、H105和W111等残基可能参与关键界面形成，同时W111插入的HBc疏水口袋具有被小分子竞争占据的潜力。本研究围绕这一结构模型，建立非感染性重组蛋白体外检测体系，用于定量验证HBc--preS相互作用及其受突变体和小分子影响的规律。

本研究首先比较了三种ELISA构型，包括HBc直接包被、StrepXT捕获preS以及抗core Fab捕获HBc。结果显示，不同构型虽然分别改善了核心蛋白或preS的呈递方式，但均受到非特异背景、固相吸附、洗涤导致弱相互作用解离以及统计窗口不足等限制，难以作为后续系统筛选和定量比较的主要平台。该结果提示HBc--preS相互作用更适合在减少洗涤和固定化干扰的体系中检测。

在此基础上，本研究建立了NanoBiT发光互补体系，将SmBiT融合至HBVCp182，将LgBiT融合至preS 87--114片段或相应阴性对照。初步实验表明，HBVCp182-SmBiT与LgBiT-preS组合能够产生高于阴性对照的发光信号；固定HBVCp182-SmBiT并滴定LgBiT-preS后，实验组与阴性对照之间出现可量化分离，部分条件下Z-factor超过0.5，说明该体系能够报告HBc--preS相互作用并具备进一步优化空间。后续通过重新纯化HBVCp182-SmBiT、比较Q柱合并组分活性以及优化LgBiT构建体，进一步提高了实验体系的可用性。

小分子抑制实验表明，在含foldon体系中，Triton X-100、NP-40和香叶醇直接加入时抑制趋势不够稳定；经预孵育并提高小分子浓度后，高浓度区间可观察到更明显的发光信号下降。然而foldon三聚化和GST二聚化可能增强表观结合，限制剂量--反应曲线的定量解释。进一步去除foldon后，Triton X-100和香叶醇在高浓度区间表现出更清楚的相对抑制趋势，但高浓度小分子对阴性对照信号的影响提示仍需加入更多背景控制。最后，在去除GST标签后的体系中，preS R102A、H105A和W111A突变均显著降低HBc--preS互作信号，其中R102A和H105A下降最明显，支持这些残基在HBc识别界面中的关键作用。

\section{主要结论}

本研究得到以下主要结论。第一，HBc--preS相互作用在常规固相ELISA体系中难以获得稳定、低背景且具有足够统计窗口的读数，说明该相互作用可能具有弱亲和力、构象敏感和多价依赖等特点。第二，NanoBiT体系比ELISA更适合本课题的体外定量检测需求，能够在接近溶液状态的条件下将HBVCp182与preS 87--114之间的接近事件转化为可检测发光信号。第三，HBVCp182-SmBiT的纯度和LgBiT-preS构建体的多聚化状态会显著影响检测窗口和定量解释，因此蛋白质量控制、foldon去除和GST标签影响评估是平台优化中的关键步骤。第四，preS R102、H105和W111对HBc--preS结合具有重要贡献，实验结果与结构模型中这些残基参与静电作用、氢键或疏水嵌合的解释相一致。第五，部分小分子在优化后的NanoBiT体系中可降低HBc--preS相关发光信号，提示HBc疏水口袋具有被小分子干预的可能性，但目前数据仍不足以直接给出精确IC$_{50}$或结合常数。

\section{不足与展望}

尽管本研究已经建立了可用于机制验证的体外NanoBiT平台，仍存在若干需要进一步改进的方面。首先，当前实验主要基于阶段性重复和有限批次蛋白样品，后续需要增加独立蛋白批次和独立实验重复，以更严格地评估平台稳定性、Z-factor范围和不同构建体之间的可比性。其次，小分子抑制实验中高浓度条件可能同时影响蛋白稳定性、NanoLuc片段互补效率或体系背景，因此需要加入更完整的控制组，例如仅检测NanoBiT片段互补本身、小分子对底物发光体系的影响、蛋白聚集或变性控制等。只有在背景变化可控的浓度范围内，才适合进一步进行剂量--反应拟合和IC$_{50}$估算。

其次，preS突变体分析目前集中于R102A、H105A和W111A三个关键位点。后续可以扩展至P100、L101、S113以及HBc侧对应界面残基，从preS和HBc两个方向系统验证结构模型。若能够结合更多点突变、组合突变和互补突变实验，将有助于区分静电作用、疏水嵌合和氢键网络对整体结合的相对贡献。对于小分子研究，也可在Triton X-100和香叶醇等初步候选分子的基础上，选择结构更明确、药物样性质更好的化合物进行验证，并结合计算对接或结构类似物比较，寻找更适合后续优化的分子骨架。

最后，本研究聚焦于非感染性重组蛋白体外体系，适合在较高安全性和可控性条件下解析HBc--preS相互作用本身。后续若要进一步评价该界面对HBV生命周期的影响，需要在符合生物安全规范和实验室审批要求的前提下，选择更接近细胞内环境的验证体系，例如非感染性共表达、蛋白共定位或病毒样颗粒相关模型。通过将体外定量结果、结构模型和更复杂体系中的功能验证相结合，可以更全面地评估HBc--preS界面作为新型抗HBV药物靶点的可行性。
